\documentclass[10pt,a4paper]{amsart}
\usepackage[margin=19mm]{geometry}
\usepackage{amsmath,amssymb,mathtools}
\usepackage[scr=rsfs,cal=euler]{mathalfa}
\usepackage{xfrac}
\usepackage[unicode,hidelinks,bookmarks=false]{hyperref}
\newcommand{\ie}{i.e.\ }
\newcommand{\cf}{cf.\ }
\newcommand{\resp}{respectively\ }
\newcommand{\Subsection}{\subsection}
\providecommand{\nobreakdash}{\nobreak\hbox{-}}
\setlength{\emergencystretch}{2em}
\multlinegap0pt
\usepackage{amssymb}
\usepackage{enumerate}
\usepackage{tikz}
\newtheorem{lem}{Lemma}
\usepackage{cleveref}
\crefname{lem}{Lemma}{Lemmas}
\newcommand{\C}{{\mathbb C}}
\newcommand{\R}{{\mathbb R}}
\title{(Positive) Quadratic Determinantal Representations of Quartic Curves and the Robinson Polynomial}
\author{Clemens Br\"user, Mario Kummer}
\date{Source: arXiv:2412.02319v1 (CC BY 4.0)}
\begin{document}
\maketitle

\noindent\emph{Source and licence.} (Positive) Quadratic Determinantal Representations of Quartic Curves and the Robinson Polynomial. Version arXiv:2412.02319v1 (CC BY 4.0), distributed by arXiv under the Creative Commons Attribution 4.0 International licence, http://creativecommons.org/licenses/by/4.0/, as recorded in the arXiv licence metadata for this exact version and shown on the abstract page https://arxiv.org/abs/2412.02319v1. Full licence notice: \texttt{licenses/arxiv-2412.02319-CC-BY-4.0.txt}.\par\smallskip

\noindent\textit{Editorial scope.} Counted unit: the article's lem lem:indep (source0.tex lines 986--988). The article's complete original proof of it is delivered with it (source0.tex lines 990--1017, one \texttt{proof} environment of 28 lines). No mathematics is written, repaired or re-derived: the statement and the proof are the article's own lines, in the article's own notation, and no retained line is substituted or re-flowed.  Retained uncounted as the source-local context the proof needs: lines 1234--1243.  Nothing was omitted from any retained span, so the package carries no omission record.  No retained line contains a citation, so the wrapper carries no bibliography and no bibliography entry.  No retained line contains a figure environment, an image-inclusion command or drawing code, so no image is delivered and the package declares no asset dependency.  Editorial formatting only, in addition: an AMS \texttt{amsart} preamble that restates the article's own declarations and macros used by the retained lines, namely the environments \texttt{lem} on separate counters, together with the standard packages those lines need.  The retained environments print in this document as Lemma 1 in order of appearance, whereas the article prints the same items with the numbering of its own full document.  The complete native source of the article as distributed in the arXiv e-print for this exact version is bundled unchanged as \texttt{sources/169521/source0.tex}.
\begin{lem} \label{lem:validchoice}
    All choices in the above list are possible. Precisely we have the following:
    \begin{enumerate}
        \item All even theta characteristics on $C$ are non-vanishing. Hence, $T_0$ exists as desired.

        \item $\ell(T_0 + H) \geq 1$. Hence, $T$ exists as desired.

        \item $\ell(T + B - H) = 2$. Hence, $A$ exists as desired.
    \end{enumerate}
\end{lem}

\begin{lem} \label{lem:indep}
    The choice of $a_{10}$ in Dixon's algorithm does not affect whether the resulting quadratic representation $Q$ is positive.
\end{lem}

\begin{proof}
    First note that given $a_{00}$ and $a_{10}$ the conic $a_{11}$ is already uniquely determined. Indeed, if through Max Noether's fundamental theorem we obtained identities
    \begin{align*}
        a_{10}^2 = a_{00}a_{11} + pF \\
        a_{10}^2 = a_{00}a_{11}' + p'F,
    \end{align*}
    then necessarily
    \begin{equation*}
        a_{00}(a_{11} - a_{11}') = (p-p')F
    \end{equation*}
    with $p - p' \in \C$, which proves that either $F$ is reducible or both sides are identically zero. By assumption on $F$ only the latter can occur, so $a_{11} = a_{11}'$.

    Since $\ell(A) = 2$ by \Cref{lem:validchoice}, if we replace $a_{10}$ by another conic $a_{10}'$ we necessarily have a relation
    \begin{equation*}
        {a_{10}'} = \lambda a_{00} + \mu a_{10}
    \end{equation*}
    for some $\lambda, \mu \in \R$. Without loss of generality we assume $\lambda = 1$. This new choice gives rise to the unique new quadratic representation
    \begin{equation*}
        Q' = \begin{pmatrix}
            1 & 0 \\
            1 & \mu
        \end{pmatrix} Q \begin{pmatrix}
            1 & 1 \\
            0 & \mu
        \end{pmatrix}.
    \end{equation*}
    In particular, due to Sylvester's law of inertia, the signature of $A'$ is the same as that of $A$ in every point of evaluation, which proves independence from choice of $a_{10}$.
\end{proof}

\end{document}
